\subsection{Post-submission 8B free-generation follow-up}
\label{app:b3-generation}

This post-submission follow-up evaluates the frozen 8B hosts, selected swaps,
and published Swaps under free generation. It uses 250 MGSM-Rev2 test
questions per language and candidate, at revision
\artifact{1463c6dc}\ldots{} (test split). Selection searched all 666 windows
and chose zero-based inclusive layers $[13,15]$ for French and $[20,31]$ for
Chinese; published windows are $[13,22]$ and $[13,20]$. Exact model commits,
prompts, runs, outputs, and audits are recorded under
\artifact{results/b3/generation\_followup\_v1/}.

Selected swaps were ordinary checkpoints made by copying English-donor layers
into each host. Per-language receipts verify 399 tensor names, the exact
selected mappings, and equality of each official Swap with its reconstruction.
All candidates used the pinned host tokenizer, localized reasoning prompt,
and enabled thinking. Sampling used temperature $1$, top-$p=0.95$, top-$k=20$,
min-$p=0$, paired per-item seeds, and an 8,192-token cap. Inference used vLLM
0.11.0 with BF16 and no quantization on one A100 SXM4 80GB; package versions
and the image digest are in the execution manifest. This differs from the
32,768-token example and ten-run averages reported by the source model cards
\citep{lasbordes2026rethinking}.

A separate archived German BF16 pilot stopped because cached and uncached
greedy decoding diverged. Its runbook reports the same disagreement for the
raw model and says the observed divergence disappeared in float32; changing
precision would have violated that pilot's frozen protocol. Partial outputs
were deleted, so no token-level trace remains for reanalysis. This does not
establish a French or Chinese generation failure. All six present
materialized-checkpoint vLLM runs completed and passed their audits.

Table~\ref{tab:b3-generation} reports frozen integer EM beside post-hoc
numeric EM; the original score is unchanged. V2 parses only the completed,
non-truncated post-think final line, using the suffix after the last answer
marker or a standalone scalar, and compares decimals and fractions exactly.
Host-only QA preceded the v1 freeze. Later comparative QA found 66 answers
accepted by integer EM but rejected by v1 because bold Markdown around markers
leaked into the suffix. V2 strips only bold, underscore, and backtick tokens
before marker extraction. Specified after comparative QA, v2 was not externally
preregistered; v1 and original scores remain intact. V2 still rejects a French
published answer ending in \artifact{64 \$} under the unchanged currency
policy. Unparseable and truncated answers count as incorrect. Both raw-output
and v2 audits passed; see
\artifact{results/b3/generation\_followup\_v1/audit\_report.json} and
\artifact{results/b3/generation\_followup\_v1/scoring\_v2\_audit.json}.

Target ID is fastText top-1 on a complete reasoning trace with at least 20
letters; incomplete, too-short, or truncated traces are unknown and count as
non-target among all 250 records \citep{grave2018fasttext}. EN span flags
English in a non-overlapping chunk of up to 200 characters with at least 20
letters and English probability at least $0.5$. These are classifier outputs,
not human-validated judgments. Of the selected Chinese model's 248 complete traces, 225
are classified as English and 23 as Chinese; two are unknown. English is
flagged in 238 of 250 traces. The published French model has 55 truncations,
compared with three for selected, so that comparison is cap-sensitive and does
not establish uncapped superiority.

Paired 95\% percentile intervals use 20,000 item resamples paired by question
ID. They are conditional on one shared deterministic seed schedule, are not
adjusted across contrasts, and do not capture generation-seed variability.
Constrained first-order weight attribution selects the same windows, so this three-choice
panel cannot attribute any difference to SUTURE over that baseline. It does
not validate all-window free-generation ranking or general language
preservation.
